\section*{7. DESIGN PRINCIPLES AND DATA SECURITY}

\phantomsection
\addcontentsline{toc}{section}{7. DESIGN PRINCIPLES AND DATA SECURITY}

This chapter shows how design principles keep the module boundaries of Chapter 4 clear, and how QR codes, permissions, audit records, and personal data are protected.

\subsection*{7.1. SOLID Principles}

\phantomsection
\addcontentsline{toc}{subsection}{7.1. SOLID Principles}

The modules own distinct decisions and data. Gate Core owns factory scan and passage records, Access \& Configuration owns access rights and station settings, Attendance owns consolidation and review, and Payroll Handover owns batches. Modules exchange data through public contracts, and the two external adapters implement ports defined by Gate Core and Payroll Handover. We apply SOLID to these modules, interfaces, and ports, not only to classes.

\subsubsection*{7.1.1. Single Responsibility Principle (SRP)}

\phantomsection
\addcontentsline{toc}{subsubsection}{7.1.1. Single Responsibility Principle (SRP)}

SRP says that a unit should have one cohesive responsibility and one kind of reason to change. It does not mean that a component may perform only one operation.

\textbf{In TMC:}

\begin{itemize}[leftmargin=36pt, itemsep=2pt, topsep=4pt, parsep=2pt]
  \item Gate Transaction Processor receives a scan, handles repeated requests, asks for the eligibility check, stores the accepted attempt, and coordinates opening, passage, and closing. All of these belong to one gate transaction; accounts, shift rules, and reports live elsewhere.
  \item Access Decision only decides eligibility. A change in how the controller is called belongs in Turnstile Adapter.
  \item Under I-06, shift rules affect Work Schedule Policy and attendance aggregation, not the gate decision. If TMC also wants shifts to restrict physical access, Access Decision needs an explicit access-policy rule, still separate from the QR signature check.
  \item Attendance Review manages cases and approval authority, while Attendance Consolidation applies the right decision version, prevents duplicates, and keeps links to the original data.
\end{itemize}

\textbf{Benefit:} a change to the approval workflow or to the controller does not force a rewrite of attendance aggregation or of the gate rules.

\subsubsection*{7.1.2. Open--Closed Principle (OCP)}

\phantomsection
\addcontentsline{toc}{subsubsection}{7.1.2. Open--Closed Principle (OCP)}

OCP aims to add new capabilities while leaving stable business logic alone. In TMC, the clearest extension points are the two external integrations.

\textbf{In TMC:}

\begin{itemize}[leftmargin=36pt, itemsep=2pt, topsep=4pt, parsep=2pt]
  \item A new controller that keeps the meaning of the turnstile ports (opening and closing results, passage and lane-clear events, emergency release and reset) only needs a new adapter, and the gate rules do not change.
  \item An HR/Payroll format change that keeps the meaning of IPayrollSubmission only changes HR/Payroll Adapter.
  \item Work Schedule Policy keeps shift and aggregation rules in one place, so a change that fits its configuration model is not implemented twice in Attendance and Reporting. A rule outside that model still needs code changes and tests.
\end{itemize}

\textbf{Limit:} OCP does not mean that code never changes. Since each application is a Modular Monolith, replacing an adapter may still need a new release of that application. The aim is that integration changes do not spread to unrelated components.

\subsubsection*{7.1.3. Liskov Substitution Principle (LSP)}

\phantomsection
\addcontentsline{toc}{subsubsection}{7.1.3. Liskov Substitution Principle (LSP)}

LSP says that a replacement must keep the behavior its consumers rely on; matching method names and types is not enough.

\Needspace{6\baselineskip}
\textbf{In TMC:} a replacement adapter must keep this behavior:

\begin{itemize}[leftmargin=36pt, itemsep=2pt, topsep=4pt, parsep=2pt]
  \item For the turnstile ports, a command receipt is not OPEN\_CONFIRMED. openGate returns only OPEN\_CONFIRMED, CONTROLLER ERROR when the controller explicitly confirms it did not open, or UNDETERMINED. Passage is a later event, LANE\_CLEAR allows closeGate, and only the controller's close confirmation means CLOSED. emergencyRelease succeeds only on EMERGENCY\_RELEASED, and resetEmergency only on RESET\_CONFIRMED. Both can instead return explicit CONTROLLER ERROR or UNDETERMINED; requesting reset alone never ends the hold. A timeout after the command may have been sent stays undetermined.
  \item Gate Core records automatic attendance only when the passage event of I-01 is received and stored. A controller without reliable passage evidence cannot satisfy I-07 by returning an opening confirmation.
  \item For IPayrollSubmission, a receipt acknowledgement (I-02) only proves that HR/Payroll received that batch. When every mandatory part of the period is in CONFIRMED batches, the handover is complete (Flow C). IC-01 also needs the data to be ready for reconciliation before noon, so either the acknowledgement contract must guarantee this or HR/Payroll must give separate evidence; this has to be agreed with its owner.
  \item Status queries are used only where the API supports them. If a controller lacks the passage or safe-closing behavior assumed in I-01 and I-05, the integration is revised, and reconciliation is used where automatic attendance cannot be justified.
\end{itemize}

\textbf{Benefit:} an adapter that passes the same contract tests can replace the old one without any change in Gate Core or Payroll Handover.

\Needspace{13\baselineskip}
\subsubsection*{7.1.4. Interface Segregation Principle (ISP)}

\phantomsection
\addcontentsline{toc}{subsubsection}{7.1.4. Interface Segregation Principle (ISP)}

ISP says that a consumer should depend only on the operations it uses.

\Needspace{6\baselineskip}
\textbf{In TMC:} the grouped interfaces in Table 4.20 are split into narrower design contracts:

\begin{itemize}[leftmargin=36pt, itemsep=2pt, topsep=4pt, parsep=2pt]
  \item \textbf{Turnstile ports:} Gate Core defines narrow contracts instead of one large interface. IOpenGate (openGate with a time limit), ICloseGate (closeGate), and IEmergencyGate (emergencyRelease and reset) are commands that return their results, and Turnstile Adapter implements them against the verified controller API (I-01, I-05). IPassageEvents (PASSAGE and LANE\_CLEAR) runs the other way: Gate Core implements it, and Turnstile Adapter calls it when the controller reports these events. Gate Transaction Processor uses the normal-cycle contracts, and the emergency coordinator uses IEmergencyGate. IControllerStatusQuery is a separate optional contract, used only where the controller API supports it.
  \item \textbf{GateCoreApi:} IScanSubmission for Station Client, IInterventionRequest for Gate Intervention (support case, assisted opening, emergency release and reset), and IFactoryUpdates for Site Sync.
  \item \textbf{Other APIs:} AttendanceApi separates review and approval from period-data queries, AuthorizationApi separates account administration from authorization checks, and AuditApi separates record intake from queries.
\end{itemize}

\textbf{Benefit:} each consumer sees only the operations it needs, so a change to one contract does not affect the other consumers.

\textbf{Limit:} these are design contracts, not a security control on their own. Server-side authorization still checks every operation and data scope.

\subsubsection*{7.1.5. Dependency Inversion Principle (DIP)}

\phantomsection
\addcontentsline{toc}{subsubsection}{7.1.5. Dependency Inversion Principle (DIP)}

Business code should depend on abstractions that it defines, not on the details of external systems.

\textbf{In TMC:}

\begin{itemize}[leftmargin=36pt, itemsep=2pt, topsep=4pt, parsep=2pt]
  \item Gate Transaction Processor depends on the command ports that Gate Core defines (IOpenGate, ICloseGate, IEmergencyGate), not on the Controller API format. Turnstile Adapter reports events back through IPassageEvents, which Gate Core also defines, so both directions point to Gate Core's abstractions.
  \item Gate Core also defines its storage ports. A unit-of-work port (IUnitOfWork) owns each transaction, so reserving the lane, claiming the token, and storing ACCEPTED happen in one write without Gate Transaction Processor knowing which database is used.
  \item Payroll Handover depends on IPayrollSubmission without knowing HR/Payroll connection details or response structures.
  \item Test adapters return command outcomes and deliver PASSAGE and LANE\_CLEAR through IPassageEvents. Receipt-only and timeout cases produce UNDETERMINED. A shared contract test suite checks substitution; integration tests check these meanings against the real controller.
\end{itemize}

\textbf{Limit:} not every internal call needs an interface. We add an abstraction only when it supports substitution, testing, or isolation of a real dependency.

\Needspace{9\baselineskip}
\subsection*{7.2. Cohesion, Coupling and Extensibility}

\phantomsection
\addcontentsline{toc}{subsection}{7.2. Cohesion, Coupling and Extensibility}

\subsubsection*{7.2.1. High Cohesion}

\phantomsection
\addcontentsline{toc}{subsubsection}{7.2.1. High Cohesion}


Functions belong in the same module when they serve related responsibilities:

\begin{itemize}[leftmargin=36pt, itemsep=2pt, topsep=4pt, parsep=2pt]
  \item Gate Core concentrates on scan processing at a factory, Attendance on recording, review, and aggregation, and Payroll Handover on preparing data and tracking its delivery.
  \item Inside a module, components keep distinct responsibilities. Access \& Configuration manages several kinds of business data, but QR issuance, access rights, stations, and shift policy are separate components.
  \item Modules are not built around data entities. There is no ``Employee module'': employee data is used in several places, but access decisions, aggregation, and account administration change for different reasons.
\end{itemize}

\Needspace{8\baselineskip}
\subsubsection*{7.2.2. Low Coupling}

\phantomsection
\addcontentsline{toc}{subsubsection}{7.2.2. Low Coupling}


Coupling is kept where it is needed and removed elsewhere:

\begin{itemize}[leftmargin=36pt, itemsep=2pt, topsep=4pt, parsep=2pt]
  \item The gate flow waits synchronously only for the eligibility check and the opening confirmation. Passage recording and safe closure run independently of the kiosk response.
  \item Gate Core stores source records and outgoing events at the factory, and Site Sync moves them later, so the kiosk never waits for central consolidation.
  \item Events reduce waiting but do not make producers and consumers independent. Both sides still depend on event structures, versions, progress, and acknowledgement rules, so contract changes are checked on both sides.
  \item Deploying three applications does not by itself create three independent quanta; Chapter 3 judges quanta from the actual dependencies.
\end{itemize}

\Needspace{15\baselineskip}
\subsubsection*{7.2.3. Extensibility}

\phantomsection
\addcontentsline{toc}{subsubsection}{7.2.3. Extensibility}


Table 7.1 shows how the design handles some hypothetical changes. They are only used to test the boundaries and do not add requirements.

\begin{longtable}{|>{\raggedright\arraybackslash}p{75.5pt}|>{\raggedright\arraybackslash}p{158.5pt}|>{\raggedright\arraybackslash}p{176.5pt}|}
\hline
\multicolumn{1}{|>{\columncolor{tableheader}\centering\arraybackslash}p{75.5pt}|}{\textbf{Scenario}} & \multicolumn{1}{>{\columncolor{tableheader}\centering\arraybackslash}p{158.5pt}|}{\textbf{Areas requiring changes}} & \multicolumn{1}{>{\columncolor{tableheader}\centering\arraybackslash}p{176.5pt}|}{\textbf{Conditions to preserve}} \\
\hline
\endhead
\multicolumn{3}{c}{\rule{0pt}{20pt}\normalsize\textbf{Table 7.1. Scope of Changes in Selected Extension Scenarios}} \\
\endlastfoot
Add a factory & Deploy the factory application and database; register stations and lanes; configure access rights and Site Sync. & Bind QR tokens to the correct factory; share token-claim state across its lanes; track synchronization progress. \\
\hline
Replace a controller & Update Turnstile Adapter and connection configuration. & Preserve opening, passage, lane-clear, and safe-closing semantics; verify I-01 and I-05. \\
\hline
Change the HR/Payroll API & Update HR/Payroll Adapter and data mappings. & Preserve batch identifiers, content versions, and acknowledgement semantics; reassess I-02. \\*
\hline\noalign{\nobreak}
Adjust attendance aggregation rules & Update Work Schedule Policy and the relevant rule-application logic. & Keep rule usage consistent across Attendance and Reporting; preserve versions and adjustment history. \\
\hline
\end{longtable}
\phantomsection\addcontentsline{lot}{table}{\protect\numberline{7.1}Scope of Changes in Selected Extension Scenarios}


The boundaries show where a change starts, but impact testing is still needed. Adding a factory, for example, also means configuring event sources, access permissions, verification keys, and data-reception progress. Maintainability stays a supporting characteristic, as in Chapter 2, so we do not add services only to make the design look extensible.

\Needspace{9\baselineskip}
\subsection*{7.3. QR Authentication and Anti-replay}

\phantomsection
\addcontentsline{toc}{subsection}{7.3. QR Authentication and Anti-replay}

\subsubsection*{7.3.1. QR Issuance and Validation}

\phantomsection
\addcontentsline{toc}{subsubsection}{7.3.1. QR Issuance and Validation}


QR credentials follow ADR-11 and I-09:

\begin{itemize}[leftmargin=36pt, itemsep=2pt, topsep=4pt, parsep=2pt]
  \item \textbf{Issuance.} An authenticated employee requests a QR code through Employee Portal, and QR Credential Issuance issues a token valid for 60 seconds and restricted to one factory. The 60-second value is our design policy, not a requirement in the assignment. Under I-08, company email only links to the authenticated live display; we do not email a reusable QR image, and whether direct display in email is required is still to be confirmed.
  \item \textbf{Content.} The signed content holds the token identifier, the employee identifier, the permitted factory, and validity information. Names, phone numbers, and attendance details are left out.
  \item \textbf{Signature.} The central site keeps the private signing key, and factories receive only the public keys needed for verification; the private key is never copied to kiosks or factory applications. A signature protects integrity and identifies the issuer, but it does not encrypt the content. Algorithms and libraries are chosen during implementation; if JWT is used, aud, exp, and jti can hold the audience, expiry, and token identifier.
  \item \textbf{Validation.} Access Decision checks the token structure, issuer, signature, validity period, and factory scope, then checks the employee's status and access rights in Eligibility Snapshot. A valid signature alone is not enough if the employee is blocked, lacks permission, or the access data is no longer usable.
\end{itemize}

\Needspace{9\baselineskip}
\subsubsection*{7.3.2. Anti-replay within a Factory}

\phantomsection
\addcontentsline{toc}{subsubsection}{7.3.2. Anti-replay within a Factory}


Duplicate requests and QR replay are different problems: scanId identifies one processing attempt, while tokenId identifies one issued QR code.

\begin{itemize}[leftmargin=36pt, itemsep=2pt, topsep=4pt, parsep=2pt]
  \item \textbf{Same scanId.} If Station Client resends a scanId after a lost response, Gate Transaction Processor returns the stored status without another command or transaction.
  \item \textbf{Claimed tokenId.} A new scanId with an already claimed tokenId is rejected as a replay. A new token with a different identifier is a new scan, even if it comes soon after another.
  \item \textbf{Atomic claim.} All lanes of a factory share Gate Core's token-claim data, and the claim is stored in the same transaction as ACCEPTED. Unique constraints make sure that only one scan is accepted when several lanes receive the same token at once; a separate check-then-claim step would leave a race.
  \item \textbf{No release after uncertainty.} Once claimed, a token never becomes unused again because of a timeout, a lost response, or a restart. The original scan is queried or reconciled instead.
\end{itemize}

\Needspace{9\baselineskip}
\subsubsection*{7.3.3. Anti-replay between Factories}

\phantomsection
\addcontentsline{toc}{subsubsection}{7.3.3. Anti-replay between Factories}


Each token is tied to one factory:

\begin{itemize}[leftmargin=36pt, itemsep=2pt, topsep=4pt, parsep=2pt]
  \item A token is valid only at the factory named in its signed content. A token for Binh Duong is rejected at Long An, even if the employee may enter both; the employee requests a separate token there.
  \item The actual factory comes from server configuration and the authenticated, registered station, not from a factory field sent by the client.
  \item This avoids a global token-claim service on every scan, because each factory controls only its own tokens.
  \item Single-use tokens prevent reuse, but they do not stop an employee from sharing an unused image, and they do not prove who is holding the phone.
\end{itemize}

\subsubsection*{7.3.4. Disconnected Operation and Recovery}

\phantomsection
\addcontentsline{toc}{subsubsection}{7.3.4. Disconnected Operation and Recovery}

QR validity (60 seconds) and the age of eligibility data (normal for 60 seconds, then degraded while the last confirmed update is less than 30 minutes old) are separate limits; Section 5.7.2 gives the full conditions. QR Credential Issuance stays central, so new tokens need that application. After failover or a backup restore, Gate Core accepts scans only when the token-claim state is consistent; otherwise the factory follows the manual procedure. Cleaning up claim data must never make a token that could still be accepted look unused, so expiry, clock error, and recovery are considered before claim state is deleted. The same holds when a scan stops after ACCEPTED because its deadline passed before openGate: the lane is released, but the token stays claimed.

\Needspace{9\baselineskip}
\subsection*{7.4. Role-based and Scoped Access Control}

\phantomsection
\addcontentsline{toc}{subsection}{7.4. Role-based and Scoped Access Control}


Permission to use TMC functions is different from permission to enter a factory. Identity \& RBAC decides which operations and data a user may access, while Access Decision checks physical access rights in the QR flow. Table 7.2 summarizes the boundaries for the six user groups; exact permissions follow TMC's approved authority and job scope.

\begin{longtable}{|>{\raggedright\arraybackslash}p{75.5pt}|>{\raggedright\arraybackslash}p{190.5pt}|>{\raggedright\arraybackslash}p{145.5pt}|}
\hline
\multicolumn{1}{|>{\columncolor{tableheader}\centering\arraybackslash}p{75.5pt}|}{\textbf{User group}} & \multicolumn{1}{>{\columncolor{tableheader}\centering\arraybackslash}p{190.5pt}|}{\textbf{Functions within the assigned scope}} & \multicolumn{1}{>{\columncolor{tableheader}\centering\arraybackslash}p{145.5pt}|}{\textbf{Boundaries to preserve}} \\
\hline
\endhead
\multicolumn{3}{c}{\rule{0pt}{20pt}\normalsize\textbf{Table 7.2. Permission Boundaries by User Group}} \\
\endlastfoot
Employee & Obtain QR codes, view personal data, and submit attendance concerns. & Cannot view other employees' data or approve their own adjustments. \\
\hline
Production Supervisor & View and monitor data within the assigned area of responsibility. & Does not receive default permission to modify attendance data directly. \\
\hline
HR Staff & Manage business data, review cases, and approve decisions within their authority. & Support cases without a QR code require HR approval under I-04. \\
\hline
Security Guard & Create support cases and perform authorized gate operations. & Cannot independently approve attendance recording for support cases. \\
\hline
System Administrator & Perform technical administration and manage accounts and configuration according to assigned permissions. & Does not receive default permission to read or modify business data. \\*
\hline\noalign{\nobreak}
Executive Management & View authorized reports and aggregated information. & Detailed personal data is available only with the corresponding permission. \\
\hline
\end{longtable}
\phantomsection\addcontentsline{lot}{table}{\protect\numberline{7.2}Permission Boundaries by User Group}


\Needspace{6\baselineskip}
The following rules apply to all groups:

\begin{itemize}[leftmargin=36pt, itemsep=2pt, topsep=4pt, parsep=2pt]
  \item \textbf{Server-side checks.} Authorization is enforced at the server, not by hiding buttons. The server checks the action, the target, and the scope, such as the factory, the assigned unit, or the user's own records. Viewing a report does not include permission to export all employee data.
  \item \textbf{Device identity.} Station Client may submit and query scans only for its registered station and cannot call approval or administration functions. A station identity and an employee's QR token do not replace each other.
  \item \textbf{No login per scan.} The QR flow does not need a central user login for each scan: the station is authenticated as a device, and eligibility is checked at the factory. Software-based assistance and emergency operations still need authorization, and the check is never bypassed when it cannot be completed.
\end{itemize}

\Needspace{9\baselineskip}
\subsection*{7.5. Audit Logging and Personal Data Protection}

\phantomsection
\addcontentsline{toc}{subsection}{7.5. Audit Logging and Personal Data Protection}

\subsubsection*{7.5.1. Audit Content}

\phantomsection
\addcontentsline{toc}{subsubsection}{7.5.1. Audit Content}


Audit records answer who did what, to which object, when, and with what result:

\begin{itemize}[leftmargin=36pt, itemsep=2pt, topsep=4pt, parsep=2pt]
  \item \textbf{What is recorded.} Authentication, access denials, gate results, assistance and emergency operations, access-right changes, adjustment approvals, and HR/Payroll handover.
  \item \textbf{Record fields.} The acting user or device, the action, the target, the time, the outcome, and correlation identifiers: scanId, commandId, and eventId for scans, and the decision or batch identifier with its version for cases and batches.
  \item \textbf{Event time.} The actual event time is kept apart from the time of later data entry, so a record entered after an incident does not suggest that the entry or exit happened later.
  \item \textbf{What is never logged.} Complete QR tokens, passwords, signing keys, and authentication secrets. Personal data and case details appear only when needed for traceability, and free-text fields are escaped so that they cannot break the log structure.
\end{itemize}

\Needspace{9\baselineskip}
\subsubsection*{7.5.2. Durable Audit and Business History}

\phantomsection
\addcontentsline{toc}{subsubsection}{7.5.2. Durable Audit and Business History}


Audit evidence is kept as durably as the business data it describes:

\begin{itemize}[leftmargin=36pt, itemsep=2pt, topsep=4pt, parsep=2pt]
  \item \textbf{Evidence at the source.} Audit evidence is committed with the factory business change it describes. The recording order and write-failure handling are defined in Flow A; central delivery is not the point at which this evidence becomes durable.
  \item \textbf{Asynchronous delivery.} Audit records reach Audit Log asynchronously. During an interrupted transfer, the source retains the evidence; scans continue only while local durability and the operating conditions of Chapter 5 hold.
  \item \textbf{No double records.} A resent event keeps its identifier, so the consumer does not create a second audit record. A broker acknowledgement does not prove that the business change was applied.
  \item \textbf{History is kept.} Adjustment history keeps the original data, the approved decision, and the application result. Reconciliation never overwrites original evidence, and the trail separates an approved case from a decision that has actually been applied.
\end{itemize}

\Needspace{9\baselineskip}
\subsubsection*{7.5.3. Audit Access and Integrity}

\phantomsection
\addcontentsline{toc}{subsubsection}{7.5.3. Audit Access and Integrity}


Access to audit data is controlled and itself traceable:

\begin{itemize}[leftmargin=36pt, itemsep=2pt, topsep=4pt, parsep=2pt]
  \item Reporting checks permissions through Identity \& RBAC before calling AuditApi.query, and queries are limited to the authorized scope. Audit Log does not call Identity \& RBAC back, which avoids a cyclic dependency.
  \item Audit data is reachable only through this path; users cannot query the database directly or bypass Reporting. Permission to write audit records does not include permission to read the whole history.
  \item Storage permissions, restricted accounts, and change detection protect records against modification or deletion, and access to, administration of, and disposal of audit data are recorded as well.
  \item Protection does not mean keeping data forever. Expired data is removed under an approved policy, and technical administrators get no default right to delete business history.
  \item Monitoring data such as latency, errors, and status helps to detect incidents but does not replace evidence of who approved what.
\end{itemize}

\Needspace{9\baselineskip}
\subsubsection*{7.5.4. Personal Data Protection}

\phantomsection
\addcontentsline{toc}{subsubsection}{7.5.4. Personal Data Protection}


For deployment in Vietnam, we use Law No. 91/2025/QH15 on Personal Data Protection and Decree No. 356/2025/ND-CP as the legal baseline. Article 3(2)--(4) of the Law covers defined purposes, suitable retention, and safeguards, and Article 25(2)--(3) covers employee-data retention and informing employees about the technical measures used to manage them.\footnote{Official texts: \href{https://congbaocdn.chinhphu.vn/CongBaoCP/CongBao/2025/7/45574/57723-1-971-972.pdf}{Law No. 91/2025/QH15}; \href{https://vanban.chinhphu.vn/?classid=1\&docid=216387\&pageid=27160}{Decree No. 356/2025/ND-CP}.} The design applies them as follows:

\begin{itemize}[leftmargin=36pt, itemsep=2pt, topsep=4pt, parsep=2pt]
  \item \textbf{Collect only what is needed.} TMC collects and transfers only the data needed for access control, attendance, reconciliation, and handover. Eligibility Snapshot does not hold the full HR record, and a QR token does not carry data that can be looked up by employee identifier.
  \item \textbf{Show only what is needed.} Kiosks show only the scan result, not attendance details or HR reasons. Employees see their own records, management reports use aggregates where individual detail is not needed, and export is a separate permission.
  \item \textbf{Protect every copy.} Connections between clients, applications, the broker, and external systems are protected, and databases, replicas, outboxes, and backups are inside the protection scope. Keys and credentials never appear in source code, diagrams, or logs, and application accounts get only the permissions they need.
  \item \textbf{Identifiers are still personal data.} Using an employee identifier instead of a name reduces what is visible, but the data can still be linked to a person, and a signature does not hide the token content.
  \item \textbf{Retention is TMC's decision.} TMC must approve retention, access, and deletion rules for attendance, assistance, audit, and backup data before operation. We do not impose one common retention period without a basis.
\end{itemize}

